% !TEX root = ../../../main.tex
\section{Orthogonality and coordinate inner products}
\label{sec:linear-algebra-i:orthogonality-and-coordinate-products}

% Source: October 22 handwritten PDF, p. 6; exercise request on p. 4.
\begin{definition}
  Vectors \(v,w\) are orthogonal if \((v,w)=0\). A family
  \(\{v_j\}\) is mutually orthogonal if
  \[
    (v_i,v_j)=0\qquad\text{whenever }i\ne j.
  \]
  It is orthonormal if it is mutually orthogonal and
  \((v_j,v_j)=1\) for every \(j\).
\end{definition}

The Pythagorean identity can be expressed using only the properties of
an inner product.
% Source: October 22 handwritten PDF, p. 4, Pythagorean proof request.
\begin{exercise}
  \label{ex:lec08:pythagorean-theorem}
  Let \(v,w\) be orthogonal vectors in a real inner-product space.
  Using only the properties of the inner product, prove the Pythagorean
  identity
  \[
    \lVert v+w\rVert^2=\lVert v\rVert^2+\lVert w\rVert^2.
  \]
  \solutionlink{lec08:pythagorean-theorem}
\end{exercise}


\subsection{An inner product chosen from a basis}

A vector space does not come with a specified inner product. Fix an
ordered basis \(E=[e_1,\ldots,e_n]\). If \(v=Ea\) and \(w=Eb\),
define
\[
  (v,w)_E=a\cdot b=\sum_{j=1}^n a^jb^j.
\]
This makes \(E\) an orthonormal basis for the chosen inner product,
\[
  (e_i,e_j)_E=\delta_{ij}.
\]
Here \(\delta_{ij}\) is the \emph{Kronecker delta}, defined by
\[
  \delta_{ij}=
  \begin{cases}
    1&\text{if }i=j,\\
    0&\text{if }i\ne j.
  \end{cases}
\]
The diagonal values \((e_i,e_i)_E=1\) say that the basis vectors
have unit norm, while the off-diagonal values \((e_i,e_j)_E=0\)
say that distinct basis vectors are orthogonal. Thus the single formula
\((e_i,e_j)_E=\delta_{ij}\) records both parts of orthonormality.
The matrix of these pairings is the identity matrix \(I_n\), with the
same \(0\)--\(1\) pattern denoted by \(\delta^i_j\) in
\autoref{sec:linear-algebra-i:characteristic-polynomials-and-similarity}.

% Historical source, explicitly requested by Lucas on October 29, 2026.
% Kronecker's collected-works reprint, pp. 145 and 150, was checked directly.
\marginnote{Leopold Kronecker's paper on bilinear forms was read to
  the Berlin Academy on 15 October 1866. He describes discussions with
  Weierstrass about theta functions as the motivation for his algebraic
  investigation. In the paper, the delta expresses that a linear
  substitution followed by its inverse gives the identity. The same
  identity pattern appears here in orthonormal coordinates.
  \cite[145, 150]{Kronecker1895BilinearForms}}

Bilinearity makes the role of the delta explicit. For \(v=Ea\) and
\(w=Eb\),
\[
  \begin{aligned}
    (v,w)_E
      &=\sum_{i=1}^n\sum_{j=1}^n a^ib^j(e_i,e_j)_E\\
      &=\sum_{i=1}^n\sum_{j=1}^n a^ib^j\delta_{ij}\\
      &=\sum_{i=1}^n a^ib^i.
  \end{aligned}
\]
Only terms with equal indices survive; all cross terms vanish.
Pairing with a single basis vector therefore reads off a coordinate,
\[
  (v,e_i)_E=\sum_{j=1}^n a^j\delta_{ji}=a^i.
\]
Conversely, prescribing \((e_i,e_j)=\delta_{ij}\) determines the inner
product by bilinearity. Choosing a basis therefore provides a way to
equip the vector space with an inner product.
% Author review: p. 6 claims infinitely many inner products without
% excluding the zero space. Retain the supplied construction, not the
% unqualified count. It also tentatively says that every inner product
% has a basis giving the coordinate dot product, followed by "(ish?)".
% Orthonormal-basis existence and its proof remain deferred for clarification.
